\chapter{Implementation and Testing}
\label{chap:implementation-testing}

\section{Implementation}
\label{sec:implementation}

\subsection{Project Structure}
\label{sec:project-structure}

The project is one Git repository with three parts:
\begin{itemize}[leftmargin=1.8em,itemsep=2pt]
  \item \texttt{src/backend} -- the Spring Boot service: \texttt{controller},
        \texttt{service}, \texttt{client}, \texttt{repository}, \texttt{entity},
        \texttt{dto}, \texttt{config} and \texttt{common} packages, Flyway
        migrations in \texttt{resources/db/migration} and the tests.
  \item \texttt{src/frontend-mobile} -- the Expo app: routes in
        \texttt{src/app}, features in \texttt{src/features}, shared
        infrastructure in \texttt{src/core}, UI components in
        \texttt{src/components} and images exported from Figma in
        \texttt{assets}.
  \item \texttt{report} -- this report, and \texttt{.github/workflows} -- the
        CI workflows.
\end{itemize}

\subsection{Key Implementation Details}
\label{sec:key-implementation}

\paragraph{Session engine.}
The session store (Zustand) starts a session with a client-generated UUID and
stores the start time and the planned end time in AsyncStorage. The main
layout calls \texttt{tick()} every second and again whenever the app returns
to the foreground; the remaining time is always \textit{end time minus now},
so background time is counted correctly (FR 2.4). The store also counts music
seconds, opened articles, watched videos and played games for the summary.
When a session ends, the record is pushed to an upload queue and the queue is
flushed at once, at start-up and after every sign-in; the backend's unique key
on \texttt{(user\_id, client\_id)} makes repeated uploads harmless (FR 7.2).

\paragraph{Music player.}
One \texttt{expo-audio} player is created for the whole app, outside React,
so music continues across screens. The audio mode allows playback in silent
mode and in the background and takes exclusive audio focus, which iOS requires
before it shows the lock-screen and Control Center controls. The first track
of a session registers the player for the lock screen with the title, artist
and artwork; later tracks only update the metadata, and the controls are
removed when the session ends (FR 4.4). A failed track is skipped
automatically, at most once per track of the queue (FR 4.7). The seek bar
follows the finger while dragging and seeks only on release, so the stream is
not reloaded on every movement.

\paragraph{Downloads.}
Tracks whose artist allows downloads carry a Jamendo download URL. The app
downloads the MP3 file with \texttt{expo-file-system} into the document folder
(\texttt{music/<track id>.mp3}) and keeps the list of downloads in AsyncStorage.
When a downloaded track is played, the local file is used instead of the
stream, so the \textit{Downloaded} station works offline. \textit{Save to
phone} uses the share sheet on iOS (with \texttt{UIFileSharingEnabled} the
files also appear in the Files app) and a folder picker on Android (FR 4.6).

\paragraph{Media focus.}
A small coordinator registers the music player. Entering the TV pauses the
music only if it is playing and remembers that it did; leaving the TV resumes
the music only in that case, so music paused by the user stays paused
(FR 4.5).

\paragraph{News reader.}
The backend reads the configured VnExpress RSS feeds, parses title, summary,
image, link and date, merges and sorts them and caches the result for five
minutes. The app opens an article in an in-app browser; when the page fails
to load it offers \textit{Try again}, \textit{Open in browser} (the system
browser through \texttt{expo-web-browser}) and \textit{Back to News}
(FR 5.4, FR 5.5).

\paragraph{YouTube guard rails.}
The TV loads \texttt{m.youtube.com} in \texttt{react-native-webview} without
changing the page. A navigation filter keeps the browser on YouTube hosts,
opens other sites in the system browser and blocks app-switch links such as
\texttt{intent://}. Because YouTube changes pages with
\texttt{history.pushState}, an injected script reports the current address
after each change; the screen extracts the video ID and counts distinct IDs.
At ten videos, or when the break ends, the browser is unmounted and the end
card is shown. Playback is paused when the app goes to the background
(FR 6.1--6.3). On the web build, where an in-app browser is not available, the
TV falls back to a curated list played with the official IFrame player.

\paragraph{Game Corner.}
The games are rows of \texttt{content\_items} with the type \texttt{GAME}.
\texttt{GET /games} returns them with \texttt{playedToday} and the per-game
streak for the user's time zone. The game screen opens the page in an in-app
browser; on the first successful load it records the play locally, counts the
game for the current break and sends \texttt{POST /games/\{id\}/plays} with the
local date. The backend accepts only dates within one day of its UTC date and
inserts with \texttt{ON CONFLICT DO NOTHING}, so concurrent or repeated
requests create at most one row (FR 8.1--8.3).

\paragraph{Content proxies on the backend.}
\texttt{JamendoClient} requests tracks by tag with the MP3 stream and download
formats. Jamendo sometimes answers \textit{success} with an empty list, so the
client retries up to three times, and an empty list is never cached. The TTL
cache returns the last good list when a provider fails; when no list exists,
or when no Jamendo client ID is configured, a small set of demo tracks keeps
the Radio usable.

\paragraph{Authentication.}
\texttt{AuthService} hashes passwords with BCrypt, issues an HS256-signed JWT
access token (15 minutes) and a random refresh token (30 days) whose SHA-256
hash is stored. A refresh revokes the old token and issues a new pair. In the
app, the axios client adds the access token to every request and, on a 401
response, runs a single shared refresh and repeats the request.

\section{Deployment Details}
\label{sec:deployment-details}

\paragraph{Backend configuration.}
All secrets come from environment variables. For local development the
backend also reads an optional, git-ignored \texttt{.env} file next to
\texttt{pom.xml}; a \texttt{.env.example} file lists the keys. Real
environment variables take precedence over the file.

\begingroup
\small
\renewcommand{\arraystretch}{1.2}
\begin{longtable}{|>{\raggedright\arraybackslash}p{4.6cm}|>{\raggedright\arraybackslash}p{9.6cm}|}
\caption{Backend environment variables}\label{tab:env-vars}\\ \hline
\textbf{Variable} & \textbf{Purpose} \\ \hline
\endfirsthead
\multicolumn{2}{c}{{\bfseries \tablename\ \thetable{} -- continued from previous page}} \\ \hline
\textbf{Variable} & \textbf{Purpose} \\ \hline
\endhead
\texttt{SPRING\_PROFILES\_ACTIVE} & \texttt{dev} (default) or \texttt{prod} \\ \hline
\texttt{DB\_URL}, \texttt{DB\_USERNAME}, \texttt{DB\_PASSWORD} & PostgreSQL connection in the \texttt{prod} profile (the \texttt{dev} profile defaults to a local \texttt{breakly\_db}) \\ \hline
\texttt{JWT\_SECRET} & Key for signing access tokens; required in \texttt{prod} \\ \hline
\texttt{JWT\_ACCESS\_TTL}, \texttt{JWT\_REFRESH\_TTL} & Token lifetimes (15 minutes and 30 days by default) \\ \hline
\texttt{JAMENDO\_CLIENT\_ID} & Jamendo API client ID; without it every station plays the demo tracks \\ \hline
\texttt{YOUTUBE\_API\_KEY} & Optional key for video metadata of the curated web fallback list \\ \hline
\texttt{CORS\_ALLOWED\_ORIGINS} & Allowed origins for the web build \\ \hline
\end{longtable}
\endgroup

\paragraph{Running the system.}
\begin{enumerate}[leftmargin=1.8em,itemsep=2pt]
  \item Start PostgreSQL in Docker (container \texttt{breakly-postgres},
        database \texttt{breakly\_db}).
  \item In \texttt{src/backend}, run \texttt{./mvnw spring-boot:run}; Flyway
        applies the migrations on start-up and the API listens on port 8080.
  \item In \texttt{src/frontend-mobile}, run \texttt{npm install} and then
        \texttt{npx expo run:ios} for a development build on an iPhone or the
        simulator. Expo Go (\texttt{npx expo start}) can be used for quick UI
        checks, but background audio, the lock-screen controls and file
        sharing need the development build. The app finds the backend on the
        same Wi-Fi through the Metro host address, or through
        \texttt{EXPO\_PUBLIC\_API\_URL}.
\end{enumerate}

\paragraph{Container image and CI.}
The backend has a multi-stage Dockerfile: a Maven image with JDK 21 builds the
jar, and a JRE 21 Alpine image runs it as a non-root user. Three GitHub
Actions workflows exist: \textit{Backend CI} builds the backend with Maven on
every change under \texttt{src/backend}, \textit{Frontend CI} installs the app
dependencies and runs the TypeScript compiler, and \textit{Deploy} is prepared
for a server deployment from \texttt{main} but does not yet push to a server.

\section{Testing}
\label{sec:testing}

\subsection{Automated Tests}
\label{sec:automated-tests}

Backend tests use JUnit 5 and Spring Boot Test. Integration tests start a real
PostgreSQL database with Testcontainers, apply all Flyway migrations and call
the API through MockMvc with real JWT tokens. All tests passed on 19 November
2026 (Table~\ref{tab:automated-tests}). The app has no automated UI tests yet;
every change is checked with the TypeScript compiler (\texttt{tsc --noEmit})
and ESLint (\texttt{expo lint}), which reported no errors for the final build.

\begingroup
\small
\renewcommand{\arraystretch}{1.2}
\begin{longtable}{|>{\raggedright\arraybackslash}p{4.9cm}|>{\raggedright\arraybackslash}p{7.2cm}|>{\centering\arraybackslash}p{1.8cm}|}
\caption{Automated backend tests}\label{tab:automated-tests}\\ \hline
\textbf{Test} & \textbf{What it checks} & \textbf{Result} \\ \hline
\endfirsthead
\multicolumn{3}{c}{{\bfseries \tablename\ \thetable{} -- continued from previous page}} \\ \hline
\textbf{Test} & \textbf{What it checks} & \textbf{Result} \\ \hline
\endhead
\testname{contextLoads} & The application starts with the database and all migrations & Pass \\ \hline
\testname{listsSeededGamesForSignedInUser} & \texttt{GET /games} returns the four active games in order, with Pips last & Pass \\ \hline
\testname{requiresAuthentication} & The games endpoint rejects requests without a token (401) & Pass \\ \hline
\testname{recordingTwiceOnSameDayIsIdempotent} & The first play returns 201, a repeat on the same day 200, and only one row exists & Pass \\ \hline
\testname{rejectsDatesFarFromToday} & A play date far from today is rejected (400) & Pass \\ \hline
\testname{unknownGameReturns404} & A play for an unknown game returns 404 & Pass \\ \hline
\testname{computesCurrentAndLongestStreaks} & Current and longest streak rules, including a streak that ended yesterday & Pass \\ \hline
\testname{acceptsOnlyPlausibleLocalPlayDates} & Local play dates are accepted only within one day of the UTC date & Pass \\ \hline
\testname{parsesVnExpressRssItems} & The RSS parser reads title, link, image and date from a VnExpress item & Pass \\ \hline
\testname{ttlCacheReturnsStaleValueWhenUpstreamFails} & The cache serves the last good value when a provider fails & Pass \\ \hline
\testname{hashesRefreshTokensDeterministically} & Refresh tokens are stored as stable SHA-256 hashes & Pass \\ \hline
\end{longtable}
\endgroup

\subsection{Acceptance Test Plan}
\label{sec:acceptance-tests}

The use cases are verified manually on the reference iPhone with a
development build. Table~\ref{tab:acceptance-tests} lists the test cases; the
result column is filled in when the test is executed on the device.

\begingroup
\small
\renewcommand{\arraystretch}{1.2}
\begin{longtable}{|>{\raggedright\arraybackslash}p{1.1cm}|>{\raggedright\arraybackslash}p{1.3cm}|>{\raggedright\arraybackslash}p{5.6cm}|>{\raggedright\arraybackslash}p{4.6cm}|>{\centering\arraybackslash}p{1.3cm}|}
\caption{Acceptance test cases}\label{tab:acceptance-tests}\\ \hline
\textbf{ID} & \textbf{UC} & \textbf{Steps} & \textbf{Expected result} & \textbf{Result} \\ \hline
\endfirsthead
\multicolumn{5}{c}{{\bfseries \tablename\ \thetable{} -- continued from previous page}} \\ \hline
\textbf{ID} & \textbf{UC} & \textbf{Steps} & \textbf{Expected result} & \textbf{Result} \\ \hline
\endhead
TC-01 & UC-01 & Sign up with a new email and an 8-character password & Home Space shown; account created & \\ \hline
TC-02 & UC-01 & Sign up again with the same email & ``already registered'' error; no second account & \\ \hline
TC-03 & UC-02 & Sign in, close and reopen the app & Still signed in & \\ \hline
TC-04 & UC-03 & Sign out in Settings & Settings shows \textit{Sign in / Create account} & \\ \hline
TC-05 & UC-04 & Select 10 minutes and tap the radio & Timer starts at 10:00 and the Radio opens & \\ \hline
TC-06 & UC-04 & Lock the phone for 3 minutes during a 10-minute break & Timer shows 7:00 $\pm$ 2 s after unlocking & \\ \hline
TC-07 & UC-04 & Let a break run to 00:00 & Toast at 1:00; media stops; summary shown & \\ \hline
TC-08 & UC-05 & Select \textit{End break} and confirm & Summary shown; history row says ended early & \\ \hline
TC-09 & UC-06 & Switch the theme and restart the app & The new theme is kept & \\ \hline
TC-10 & UC-07 & Play a station, open News, then lock the phone & Music continues; lock screen shows the track & \\ \hline
TC-11 & UC-07 & Play music, open the TV, leave the TV & Music pauses and then resumes & \\ \hline
TC-12 & UC-08 & Download a track, enable airplane mode, play \textit{Downloaded} & The track plays offline & \\ \hline
TC-13 & UC-09 & Open News, filter by a category, open an article & Filtered list; article opens with the timer & \\ \hline
TC-14 & UC-10 & Open an article without a connection, then \textit{Open in browser} & Failure message; system browser opens & \\ \hline
TC-15 & UC-11 & Open ten different videos on the TV & End card after the tenth video & \\ \hline
TC-16 & UC-12 & Open Wordle, go back to the Game Corner & Wordle marked as played today & \\ \hline
TC-17 & UC-13 & Finish breaks on two consecutive days & Streak shows 2 & \\ \hline
TC-18 & UC-14 & Turn off the 1-minute warning and run a break to the end & No toast in the final minute & \\ \hline
\end{longtable}
\endgroup
